% Options for packages loaded elsewhere
\PassOptionsToPackage{unicode}{hyperref}
\PassOptionsToPackage{hyphens}{url}
\PassOptionsToPackage{dvipsnames,svgnames,x11names}{xcolor}
\documentclass[
  11pt,
  a4paper,
]{article}
\usepackage{xcolor}
\usepackage[margin=1in]{geometry}
\usepackage{amsmath,amssymb}
\setcounter{secnumdepth}{-\maxdimen} % remove section numbering
\usepackage{iftex}
\ifPDFTeX
  \usepackage[T1]{fontenc}
  \usepackage[utf8]{inputenc}
  \usepackage{textcomp} % provide euro and other symbols
\else % if luatex or xetex
  \usepackage{unicode-math} % this also loads fontspec
  \defaultfontfeatures{Scale=MatchLowercase}
  \defaultfontfeatures[\rmfamily]{Ligatures=TeX,Scale=1}
\fi
\usepackage{lmodern}
\ifPDFTeX\else
  % xetex/luatex font selection
  \ifLuaTeX
    \usepackage{luaotfload}
    \directlua{luaotfload.add_fallback("mainfontfallback",{
      "NotoSerifBengali:mode=harf"
    })}
  \fi
  \setmainfont[,RawFeature={fallback=mainfontfallback}]{Latin Modern Roman}
  \ifLuaTeX
    \usepackage{luaotfload}
    \directlua{luaotfload.add_fallback("sansfontfallback",{
      "NotoSansBengali:mode=harf"
    })}
  \fi
  \setsansfont[,RawFeature={fallback=sansfontfallback}]{Latin Modern Sans}
  \ifLuaTeX
    \usepackage{luaotfload}
    \directlua{luaotfload.add_fallback("monofontfallback",{
      "NotoSansBengali:mode=harf"
    })}
  \fi
  \setmonofont[,RawFeature={fallback=monofontfallback}]{DejaVu Sans Mono}
\fi
% Use upquote if available, for straight quotes in verbatim environments
\IfFileExists{upquote.sty}{\usepackage{upquote}}{}
\IfFileExists{microtype.sty}{% use microtype if available
  \usepackage[]{microtype}
  \UseMicrotypeSet[protrusion]{basicmath} % disable protrusion for tt fonts
}{}
\makeatletter
\@ifundefined{KOMAClassName}{% if non-KOMA class
  \IfFileExists{parskip.sty}{%
    \usepackage{parskip}
  }{% else
    \setlength{\parindent}{0pt}
    \setlength{\parskip}{6pt plus 2pt minus 1pt}}
}{% if KOMA class
  \KOMAoptions{parskip=half}}
\makeatother
\usepackage{color}
\usepackage{fancyvrb}
\newcommand{\VerbBar}{|}
\newcommand{\VERB}{\Verb[commandchars=\\\{\}]}
\DefineVerbatimEnvironment{Highlighting}{Verbatim}{commandchars=\\\{\}}
% Add ',fontsize=\small' for more characters per line
\usepackage{framed}
\definecolor{shadecolor}{RGB}{248,248,248}
\newenvironment{Shaded}{\begin{snugshade}}{\end{snugshade}}
\newcommand{\AlertTok}[1]{\textcolor[rgb]{0.94,0.16,0.16}{#1}}
\newcommand{\AnnotationTok}[1]{\textcolor[rgb]{0.56,0.35,0.01}{\textbf{\textit{#1}}}}
\newcommand{\AttributeTok}[1]{\textcolor[rgb]{0.13,0.29,0.53}{#1}}
\newcommand{\BaseNTok}[1]{\textcolor[rgb]{0.00,0.00,0.81}{#1}}
\newcommand{\BuiltInTok}[1]{#1}
\newcommand{\CharTok}[1]{\textcolor[rgb]{0.31,0.60,0.02}{#1}}
\newcommand{\CommentTok}[1]{\textcolor[rgb]{0.56,0.35,0.01}{\textit{#1}}}
\newcommand{\CommentVarTok}[1]{\textcolor[rgb]{0.56,0.35,0.01}{\textbf{\textit{#1}}}}
\newcommand{\ConstantTok}[1]{\textcolor[rgb]{0.56,0.35,0.01}{#1}}
\newcommand{\ControlFlowTok}[1]{\textcolor[rgb]{0.13,0.29,0.53}{\textbf{#1}}}
\newcommand{\DataTypeTok}[1]{\textcolor[rgb]{0.13,0.29,0.53}{#1}}
\newcommand{\DecValTok}[1]{\textcolor[rgb]{0.00,0.00,0.81}{#1}}
\newcommand{\DocumentationTok}[1]{\textcolor[rgb]{0.56,0.35,0.01}{\textbf{\textit{#1}}}}
\newcommand{\ErrorTok}[1]{\textcolor[rgb]{0.64,0.00,0.00}{\textbf{#1}}}
\newcommand{\ExtensionTok}[1]{#1}
\newcommand{\FloatTok}[1]{\textcolor[rgb]{0.00,0.00,0.81}{#1}}
\newcommand{\FunctionTok}[1]{\textcolor[rgb]{0.13,0.29,0.53}{\textbf{#1}}}
\newcommand{\ImportTok}[1]{#1}
\newcommand{\InformationTok}[1]{\textcolor[rgb]{0.56,0.35,0.01}{\textbf{\textit{#1}}}}
\newcommand{\KeywordTok}[1]{\textcolor[rgb]{0.13,0.29,0.53}{\textbf{#1}}}
\newcommand{\NormalTok}[1]{#1}
\newcommand{\OperatorTok}[1]{\textcolor[rgb]{0.81,0.36,0.00}{\textbf{#1}}}
\newcommand{\OtherTok}[1]{\textcolor[rgb]{0.56,0.35,0.01}{#1}}
\newcommand{\PreprocessorTok}[1]{\textcolor[rgb]{0.56,0.35,0.01}{\textit{#1}}}
\newcommand{\RegionMarkerTok}[1]{#1}
\newcommand{\SpecialCharTok}[1]{\textcolor[rgb]{0.81,0.36,0.00}{\textbf{#1}}}
\newcommand{\SpecialStringTok}[1]{\textcolor[rgb]{0.31,0.60,0.02}{#1}}
\newcommand{\StringTok}[1]{\textcolor[rgb]{0.31,0.60,0.02}{#1}}
\newcommand{\VariableTok}[1]{\textcolor[rgb]{0.00,0.00,0.00}{#1}}
\newcommand{\VerbatimStringTok}[1]{\textcolor[rgb]{0.31,0.60,0.02}{#1}}
\newcommand{\WarningTok}[1]{\textcolor[rgb]{0.56,0.35,0.01}{\textbf{\textit{#1}}}}
\usepackage{longtable,booktabs,array}
\usepackage{calc} % for calculating minipage widths
% Correct order of tables after \paragraph or \subparagraph
\usepackage{etoolbox}
\makeatletter
\patchcmd\longtable{\par}{\if@noskipsec\mbox{}\fi\par}{}{}
\makeatother
% Allow footnotes in longtable head/foot
\IfFileExists{footnotehyper.sty}{\usepackage{footnotehyper}}{\usepackage{footnote}}
\makesavenoteenv{longtable}
\setlength{\emergencystretch}{3em} % prevent overfull lines
\providecommand{\tightlist}{%
  \setlength{\itemsep}{0pt}\setlength{\parskip}{0pt}}
% paper_preamble.tex — shared LaTeX preamble for the PCMA/SI paper set.
%
% Compiled with lualatex (NOT pdflatex) because several papers contain
% Bengali script: the CLL paper realises one cognitive state into
% English, Bengali and Banglish, so the Bengali output is EVIDENCE and
% cannot be transliterated away without destroying the point being made.
%
% Font selection is NOT done here. pandoc >= 3 has native
% `mainfontfallback` / `monofontfallback` variables that emit the
% luaotfload fallback registration in the correct position relative to
% unicode-math; doing it by hand in this header lands after unicode-math
% and breaks font loading outright. See tools/build_papers.sh.

\usepackage{amsmath}
\usepackage{amssymb}
\usepackage{mathtools}   % \xleftrightarrow etc. used in cll.md §5
\usepackage{booktabs}

% ---------------------------------------------------------------------
% Unicode symbols appearing in PROSE rather than inside $...$.
%
% The papers were written as markdown and read as markdown, so they use
% real characters — "the ω-boost", "α-child", "ΔR coherence" — instead of
% LaTeX macros. Under unicode-math those codepoints are math-mode-only
% and abort the build ("\symbballowed only in math mode"). Mapping each
% to \ensuremath keeps the markdown readable on GitHub AND compilable,
% which is the whole point of keeping one source.
% ---------------------------------------------------------------------
\usepackage{newunicodechar}

% Greek, lower case
\newunicodechar{α}{\ensuremath{\alpha}}
\newunicodechar{β}{\ensuremath{\beta}}
\newunicodechar{γ}{\ensuremath{\gamma}}
\newunicodechar{δ}{\ensuremath{\delta}}
\newunicodechar{ε}{\ensuremath{\varepsilon}}
\newunicodechar{η}{\ensuremath{\eta}}
\newunicodechar{θ}{\ensuremath{\theta}}
\newunicodechar{κ}{\ensuremath{\kappa}}
\newunicodechar{λ}{\ensuremath{\lambda}}
\newunicodechar{μ}{\ensuremath{\mu}}
\newunicodechar{π}{\ensuremath{\pi}}
\newunicodechar{ρ}{\ensuremath{\rho}}
\newunicodechar{σ}{\ensuremath{\sigma}}
\newunicodechar{τ}{\ensuremath{\tau}}
\newunicodechar{φ}{\ensuremath{\varphi}}
\newunicodechar{χ}{\ensuremath{\chi}}
\newunicodechar{ψ}{\ensuremath{\psi}}
\newunicodechar{ω}{\ensuremath{\omega}}
\newunicodechar{ν}{\ensuremath{\nu}}
\newunicodechar{ξ}{\ensuremath{\xi}}
\newunicodechar{ζ}{\ensuremath{\zeta}}
\newunicodechar{ι}{\ensuremath{\iota}}

% Greek, upper case
\newunicodechar{Δ}{\ensuremath{\Delta}}
\newunicodechar{Φ}{\ensuremath{\Phi}}
\newunicodechar{Σ}{\ensuremath{\Sigma}}
\newunicodechar{Ψ}{\ensuremath{\Psi}}
\newunicodechar{Ω}{\ensuremath{\Omega}}
\newunicodechar{Γ}{\ensuremath{\Gamma}}
\newunicodechar{Λ}{\ensuremath{\Lambda}}
\newunicodechar{Θ}{\ensuremath{\Theta}}
\newunicodechar{Π}{\ensuremath{\Pi}}

% Operators and relations
\newunicodechar{×}{\ensuremath{\times}}
\newunicodechar{·}{\ensuremath{\cdot}}
\newunicodechar{±}{\ensuremath{\pm}}
\newunicodechar{−}{\ensuremath{-}}
\newunicodechar{≈}{\ensuremath{\approx}}
\newunicodechar{≤}{\ensuremath{\leq}}
\newunicodechar{≥}{\ensuremath{\geq}}
\newunicodechar{≠}{\ensuremath{\neq}}
\newunicodechar{∈}{\ensuremath{\in}}
\newunicodechar{∂}{\ensuremath{\partial}}
\newunicodechar{∇}{\ensuremath{\nabla}}
\newunicodechar{√}{\ensuremath{\sqrt{}}}
\newunicodechar{∞}{\ensuremath{\infty}}
\newunicodechar{∝}{\ensuremath{\propto}}
\newunicodechar{∑}{\ensuremath{\sum}}
\newunicodechar{∫}{\ensuremath{\int}}
\newunicodechar{ℝ}{\ensuremath{\mathbb{R}}}
\newunicodechar{ℕ}{\ensuremath{\mathbb{N}}}
\newunicodechar{‖}{\ensuremath{\|}}
\newunicodechar{⊕}{\ensuremath{\oplus}}
\newunicodechar{⊗}{\ensuremath{\otimes}}
\newunicodechar{∀}{\ensuremath{\forall}}
\newunicodechar{∃}{\ensuremath{\exists}}
\newunicodechar{¬}{\ensuremath{\neg}}
\newunicodechar{∧}{\ensuremath{\wedge}}
\newunicodechar{∨}{\ensuremath{\vee}}

% Arrows
\newunicodechar{→}{\ensuremath{\rightarrow}}
\newunicodechar{←}{\ensuremath{\leftarrow}}
\newunicodechar{↔}{\ensuremath{\leftrightarrow}}
\newunicodechar{↑}{\ensuremath{\uparrow}}
\newunicodechar{↓}{\ensuremath{\downarrow}}
\newunicodechar{⇒}{\ensuremath{\Rightarrow}}
\newunicodechar{⇐}{\ensuremath{\Leftarrow}}
\newunicodechar{⇔}{\ensuremath{\Leftrightarrow}}

% Super/subscripts written as literal characters
\newunicodechar{²}{\ensuremath{^{2}}}
\newunicodechar{³}{\ensuremath{^{3}}}
\newunicodechar{⁴}{\ensuremath{^{4}}}
\newunicodechar{⁵}{\ensuremath{^{5}}}
\newunicodechar{⁶}{\ensuremath{^{6}}}
\newunicodechar{₀}{\ensuremath{_{0}}}
\newunicodechar{₁}{\ensuremath{_{1}}}
\newunicodechar{₂}{\ensuremath{_{2}}}
\newunicodechar{₃}{\ensuremath{_{3}}}
\newunicodechar{ᵀ}{\ensuremath{^{\mathsf{T}}}}

% Status glyphs. The papers use emoji as status markers in tables
% (implemented / failed / warning). Mapped to typographic equivalents:
% a colour emoji font in lualatex is unreliable, and the MEANING is what
% carries — a reader needs to see "this one failed", not a red cross.
\newunicodechar{✅}{\ensuremath{\checkmark}}
\newunicodechar{✓}{\ensuremath{\checkmark}}
\newunicodechar{❌}{\ensuremath{\times}}
\newunicodechar{✗}{\ensuremath{\times}}
\newunicodechar{⚠}{\textbf{!}}
\newunicodechar{🧩}{\textbf{[piece]}}
\newunicodechar{📋}{\textbf{[list]}}
\newunicodechar{🟡}{\textbf{[partial]}}
\newunicodechar{🔴}{\textbf{[dead]}}
\newunicodechar{🟢}{\textbf{[live]}}
\newunicodechar{▼}{\ensuremath{\blacktriangledown}}
\newunicodechar{►}{\ensuremath{\blacktriangleright}}
\newunicodechar{□}{\ensuremath{\square}}
\newunicodechar{█}{\rule{0.6em}{0.9ex}}
\newunicodechar{•}{\textbullet}
\newunicodechar{†}{\textdagger}
\newunicodechar{‑}{-}
\newunicodechar{‐}{-}
\newunicodechar{…}{\ldots}
\newunicodechar{।}{\textbar}   % danda: rendered by the Bengali fallback in text
\newunicodechar{⁻}{\ensuremath{^{-}}}

% Zero-width formatting codepoints. These are invisible in an editor but
% are real characters to LaTeX: an emoji written as "⚠️" is TWO
% codepoints (U+26A0 U+FE0F), so declaring only the visible one leaves a
% variation selector with no definition and the build dies on a symbol
% nobody can see. Both are mapped to nothing.
\newunicodechar{^^^^fe0f}{}    % VARIATION SELECTOR-16
\newunicodechar{^^^^200d}{}    % ZERO WIDTH JOINER
\newunicodechar{^^^^200b}{}    % ZERO WIDTH SPACE

\usepackage{longtable}
\usepackage{array}

% Comparison matrices and cost ledgers are wide. Let them shrink to the
% text block rather than overflow the margin.
\usepackage{adjustbox}

% Code blocks: wrap long lines instead of running off the page. Several
% papers quote real C++ with long identifiers.
\usepackage{fvextra}
\fvset{breaklines=true,breakanywhere=true,fontsize=\small}

\setlength{\parskip}{0.5em}
\setlength{\parindent}{0pt}

\usepackage{fancyhdr}
\pagestyle{fancy}
\fancyhf{}
\fancyfoot[C]{\thepage}
\fancyhead[L]{\small\itshape Synthetic Intelligence — research lineage}
\fancyhead[R]{\small\itshape irfan@furylogic.com}
\renewcommand{\headrulewidth}{0.2pt}
\setlength{\headheight}{14pt}
\setlength{\headsep}{18pt}
\usepackage{bookmark}
\IfFileExists{xurl.sty}{\usepackage{xurl}}{} % add URL line breaks if available
\urlstyle{same}
\hypersetup{
  pdftitle={Future Fixes: Roadmap to Complete SI},
  pdfauthor={Irfan Mahir --- Independent Researcher, Dhaka, Bangladesh --- irfan@furylogic.com},
  colorlinks=true,
  linkcolor={black},
  filecolor={Maroon},
  citecolor={Blue},
  urlcolor={[HTML]{0B5394}},
  pdfcreator={LaTeX via pandoc}}

\title{Future Fixes: Roadmap to Complete SI}
\author{Irfan Mahir --- Independent Researcher, Dhaka, Bangladesh ---
irfan@furylogic.com}
\date{}

\begin{document}
\maketitle

{
\hypersetup{linkcolor=}
\setcounter{tocdepth}{3}
\tableofcontents
}
\section{Future Fixes: Roadmap to Complete
SI}\label{future-fixes-roadmap-to-complete-si}

\textbf{Irfan Mahir} \emph{Independent Researcher, Dhaka, Bangladesh}
irfan@furylogic.com

\begin{quote}
\emph{``Nothing on this list is proven impossible. Three are
engineering. Four are research programmes. Three are hard math. All are
tractable.''}
\end{quote}

\textbf{Companion documents.} - Theory:
\texttt{research-paper/synthetic-intelligence.md} §8 (Open Problems) -
Mechanism: \texttt{research-paper/rbg.md} - Implementation:
\texttt{research-paper/failed\_but\_parts\_works/deprected/how\_we\_implement.md}

\begin{center}\rule{0.5\linewidth}{0.5pt}\end{center}

\subsection{Overview}\label{overview}

Ten problems remain between the current PCMA engine and what §8 of the
SI paper calls ``complete SI''. Each entry below states the problem, why
it matters, the concrete fix path, the files to touch, the estimated
effort, and the success criterion. Ordered by expected difficulty
(easiest first).

\subsubsection{Effort classes}\label{effort-classes}

\begin{longtable}[]{@{}
  >{\raggedright\arraybackslash}p{(\linewidth - 4\tabcolsep) * \real{0.3333}}
  >{\raggedright\arraybackslash}p{(\linewidth - 4\tabcolsep) * \real{0.3333}}
  >{\raggedright\arraybackslash}p{(\linewidth - 4\tabcolsep) * \real{0.3333}}@{}}
\toprule\noalign{}
\begin{minipage}[b]{\linewidth}\raggedright
Class
\end{minipage} & \begin{minipage}[b]{\linewidth}\raggedright
Time budget
\end{minipage} & \begin{minipage}[b]{\linewidth}\raggedright
Nature
\end{minipage} \\
\midrule\noalign{}
\endhead
\bottomrule\noalign{}
\endlastfoot
\textbf{E} --- Engineering & Days & Function-level edits, no new
theory \\
\textbf{R} --- Research & Weeks--months & New mechanism, empirical
validation required \\
\textbf{M} --- Math & Months--year+ & Genuine theorem or categorical
construction \\
\end{longtable}

\begin{center}\rule{0.5\linewidth}{0.5pt}\end{center}

\subsection{1. Auto-tuning Fission
Thresholds}\label{auto-tuning-fission-thresholds}

\textbf{✅ SHIPPED} --- homeostatic energy feedback, EMA-smoothed, k
\textgreater{} 0 guarded. See \texttt{data/phase1/results.md}. 20-seed
drift: 0.07\% RMS.

\textbf{Class: E} · \textbf{Effort: 2--3 days} · \textbf{Files:
\texttt{src/layer2.hpp}, \texttt{src/ask.cpp}}

\subsubsection{Problem}\label{problem}

\texttt{fission\_I\_thresh\ =\ 0.15},
\texttt{fission\_tau\_thresh\ =\ 0.005},
\texttt{fission\_hold\_target\ =\ 3} are hand-picked on the
insulin-pathway ontology. On a new ontology they may fire too eagerly
(cascade the graph) or too rarely (miss real contradictions).

\subsubsection{Fix}\label{fix}

Replace fixed thresholds with quantile-based thresholds computed at the
end of every N structural ticks (say N=100):

\begin{Shaded}
\begin{Highlighting}[]
\CommentTok{// At the end of structural\_tick\_once() every N ticks:}
\DataTypeTok{void}\NormalTok{ refresh\_fission\_thresholds}\OperatorTok{(}\DataTypeTok{float}\NormalTok{ I\_quantile }\OperatorTok{=} \FloatTok{0.95}\BuiltInTok{f}\OperatorTok{,}
                                \DataTypeTok{float}\NormalTok{ tau\_quantile }\OperatorTok{=} \FloatTok{0.90}\BuiltInTok{f}\OperatorTok{)}
\OperatorTok{\{}
    \BuiltInTok{std::}\NormalTok{vector}\OperatorTok{\textless{}}\DataTypeTok{float}\OperatorTok{\textgreater{}}\NormalTok{ I\_vals}\OperatorTok{,}\NormalTok{ tau\_vals}\OperatorTok{;}
    \ControlFlowTok{for} \OperatorTok{(}\DataTypeTok{int}\NormalTok{ i }\OperatorTok{=} \DecValTok{0}\OperatorTok{;}\NormalTok{ i }\OperatorTok{\textless{}}\NormalTok{ N}\OperatorTok{;} \OperatorTok{++}\NormalTok{i}\OperatorTok{)} \OperatorTok{\{}
        \DataTypeTok{float} \VariableTok{p\_} \OperatorTok{=}\NormalTok{ pos\_updates\_by\_cavity}\OperatorTok{[}\NormalTok{i}\OperatorTok{];}
        \DataTypeTok{float} \VariableTok{n\_} \OperatorTok{=}\NormalTok{ neg\_updates\_by\_cavity}\OperatorTok{[}\NormalTok{i}\OperatorTok{];}
        \DataTypeTok{float}\NormalTok{ denom }\OperatorTok{=} \VariableTok{p\_} \OperatorTok{+} \VariableTok{n\_} \OperatorTok{+} \FloatTok{1e{-}9}\BuiltInTok{f}\OperatorTok{;}
\NormalTok{        I\_vals}\OperatorTok{.}\NormalTok{push\_back}\OperatorTok{((}\VariableTok{p\_} \OperatorTok{*} \VariableTok{n\_}\OperatorTok{)} \OperatorTok{/} \OperatorTok{(}\NormalTok{denom }\OperatorTok{*}\NormalTok{ denom}\OperatorTok{));}
\NormalTok{        tau\_vals}\OperatorTok{.}\NormalTok{push\_back}\OperatorTok{(}\NormalTok{local\_tension}\OperatorTok{[}\NormalTok{i}\OperatorTok{]);}
    \OperatorTok{\}}
    \BuiltInTok{std::}\NormalTok{sort}\OperatorTok{(}\NormalTok{I\_vals}\OperatorTok{.}\NormalTok{begin}\OperatorTok{(),}\NormalTok{ I\_vals}\OperatorTok{.}\NormalTok{end}\OperatorTok{());}
    \BuiltInTok{std::}\NormalTok{sort}\OperatorTok{(}\NormalTok{tau\_vals}\OperatorTok{.}\NormalTok{begin}\OperatorTok{(),}\NormalTok{ tau\_vals}\OperatorTok{.}\NormalTok{end}\OperatorTok{());}
\NormalTok{    fission\_I\_thresh   }\OperatorTok{=}\NormalTok{ I\_vals  }\OperatorTok{[(}\DataTypeTok{size\_t}\OperatorTok{)(}\NormalTok{I\_vals}\OperatorTok{.}\NormalTok{size}\OperatorTok{()}   \OperatorTok{*}\NormalTok{ I\_quantile}\OperatorTok{)];}
\NormalTok{    fission\_tau\_thresh }\OperatorTok{=}\NormalTok{ tau\_vals}\OperatorTok{[(}\DataTypeTok{size\_t}\OperatorTok{)(}\NormalTok{tau\_vals}\OperatorTok{.}\NormalTok{size}\OperatorTok{()} \OperatorTok{*}\NormalTok{ tau\_quantile}\OperatorTok{)];}
    \CommentTok{// fission\_hold\_target unchanged — streak requirement stays semantic}
\OperatorTok{\}}
\end{Highlighting}
\end{Shaded}

\subsubsection{Success criterion}\label{success-criterion}

\begin{itemize}
\tightlist
\item
  On \texttt{insulin\_pathway}, fission count per training round remains
  7--10 (matches current).
\item
  On a \emph{new} untested ontology (say
  \texttt{data/maxwell\_equations.axioms}), fission fires on 1--3
  concepts with genuine polarity conflict, and 0 on pure-null concepts.
\item
  20-seed reproducibility numbers change by \textless{} 5\%.
\end{itemize}

\subsubsection{Fallback}\label{fallback}

Add a shell command
\texttt{fission\_thresh\_mode\ \textless{}fixed\textbar{}quantile\textgreater{}}
so both modes can be compared side-by-side.

\begin{center}\rule{0.5\linewidth}{0.5pt}\end{center}

\subsection{2. Online L6-A Discovery}\label{online-l6-a-discovery}

\textbf{✅ SHIPPED} --- in-tick call to
\texttt{discover\_analogies\_locked()} with
\texttt{only\_parent\ =\ \textless{}new\ fissioned\textgreater{}}. Same
5 bridges as post-hoc invocation. Duplicate-bridge guard active.

\textbf{Class: E} · \textbf{Effort: 3--5 days} · \textbf{Files:
\texttt{src/layer2.hpp}}

\subsubsection{Problem}\label{problem-1}

\texttt{discover\_analogies()} currently runs \emph{after} training. In
a live-running system, analogies should appear the moment two α-children
accumulate enough shared targets --- not require an operator to invoke
discovery.

\subsubsection{Fix}\label{fix-1}

Move the intersection check into \texttt{structural\_tick\_once()},
guarded by a cadence counter to avoid running it every tick:

\begin{Shaded}
\begin{Highlighting}[]
\CommentTok{// Add to ConstructiveEngine:}
\DataTypeTok{int}\NormalTok{   online\_l6a\_cadence     }\OperatorTok{=} \DecValTok{50}\OperatorTok{;}    \CommentTok{// check every 50 structural ticks}
\DataTypeTok{bool}\NormalTok{  online\_l6a\_enabled     }\OperatorTok{=} \KeywordTok{false}\OperatorTok{;} \CommentTok{// opt{-}in}
\DataTypeTok{int}\NormalTok{   l6a\_min\_shared         }\OperatorTok{=} \DecValTok{2}\OperatorTok{;}
\DataTypeTok{int}\NormalTok{   l6a\_last\_alpha\_count   }\OperatorTok{=} \DecValTok{0}\OperatorTok{;}

\CommentTok{// Inside structural\_tick\_once(), at the end:}
\ControlFlowTok{if} \OperatorTok{(}\NormalTok{online\_l6a\_enabled}
    \OperatorTok{\&\&} \OperatorTok{(}\NormalTok{rbg\_tick\_counter }\OperatorTok{\%}\NormalTok{ online\_l6a\_cadence}\OperatorTok{)} \OperatorTok{==} \DecValTok{0}\OperatorTok{)}
\OperatorTok{\{}
    \CommentTok{// Only run if the number of fissioned parents has grown since last check.}
    \DataTypeTok{int}\NormalTok{ cur\_alpha\_count }\OperatorTok{=} \OperatorTok{(}\DataTypeTok{int}\OperatorTok{)}\NormalTok{sub}\OperatorTok{{-}\textgreater{}}\NormalTok{fission\_map}\OperatorTok{.}\NormalTok{size}\OperatorTok{();}
    \ControlFlowTok{if} \OperatorTok{(}\NormalTok{cur\_alpha\_count }\OperatorTok{\textgreater{}}\NormalTok{ l6a\_last\_alpha\_count}\OperatorTok{)} \OperatorTok{\{}
        \DataTypeTok{int}\NormalTok{ newly\_bridged }\OperatorTok{=}\NormalTok{ discover\_analogies\_locked}\OperatorTok{(}\NormalTok{l6a\_min\_shared}\OperatorTok{);}
\NormalTok{        l6a\_last\_alpha\_count }\OperatorTok{=}\NormalTok{ cur\_alpha\_count}\OperatorTok{;}
        \OperatorTok{(}\DataTypeTok{void}\OperatorTok{)}\NormalTok{newly\_bridged}\OperatorTok{;}
    \OperatorTok{\}}
\OperatorTok{\}}
\end{Highlighting}
\end{Shaded}

The \texttt{\_locked} variant assumes \texttt{sub-\textgreater{}mtx} is
already held (as it is inside the structural tick). Requires refactoring
the current \texttt{discover\_analogies()} into public and internal
variants, matching the \texttt{fission()} / \texttt{fission\_locked()}
split.

\subsubsection{Success criterion}\label{success-criterion-1}

\begin{itemize}
\tightlist
\item
  Training \texttt{metformin} then \texttt{berberine} with
  \texttt{online\_l6a\_enabled\ =\ true} produces the same bridges as
  the post-hoc \texttt{discover\_analogies} call.
\item
  Bridges appear within 100 ticks of the second concept fissioning.
\item
  No performance regression on the 20-seed benchmark (structural tick
  stays under 1 ms).
\end{itemize}

\subsubsection{Risk}\label{risk}

Bridge insertion during ongoing Hebbian consolidation could destabilize
plasticity. If observed, add a ``cooldown'' after each bridge insertion
(skip Hebbian on that tunnel for 5 ticks).

\begin{center}\rule{0.5\linewidth}{0.5pt}\end{center}

\subsection{3. Language Generation Stage 2 (Syntax
Macro-Layer)}\label{language-generation-stage-2-syntax-macro-layer}

\textbf{✅ SHIPPED} --- SyntaxSlot struct with harmonic ω ratios (3:4:5
on base 15), phases at 2πk/N, Kuramoto pass 3b, temporal readout head
with narrow-Gaussian (Q=16) exclusion.
\texttt{generate\ \textless{}label\textgreater{}} shell command produces
physics-selected slot winners, empty slots as honest refusal when no
cavity phase-locks. See \texttt{data/phase2/results.md}.

\textbf{Class: R} · \textbf{Effort: 2--3 weeks} · \textbf{Files:
\texttt{src/physics.hpp}, \texttt{src/mal.hpp},
\texttt{src/nl\_translator.hpp}}

\subsubsection{Problem}\label{problem-2}

Section 5.4b of the SI paper proposes a three-stage physics-only
language pipeline. Stages 1 (resonance signature) and 3 (label lookup)
exist. Stage 2 --- a syntax macro-layer with fixed grammatical-role
frequencies --- does not.

\subsubsection{Fix}\label{fix-2}

\textbf{Step 3.1: Add syntax slots to physics.hpp.}

\begin{Shaded}
\begin{Highlighting}[]
\KeywordTok{struct}\NormalTok{ SyntaxSlot }\OperatorTok{\{}
    \DataTypeTok{float}\NormalTok{ natural\_omega}\OperatorTok{;}      \CommentTok{// fixed: ω\_S, ω\_V, ω\_O}
    \DataTypeTok{float}\NormalTok{ phase}\OperatorTok{;}              \CommentTok{// Kuramoto{-}integrated}
    \DataTypeTok{float}\NormalTok{ amplitude}\OperatorTok{;}          \CommentTok{// sum of semantically locked cavities}
    \BuiltInTok{std::}\NormalTok{string role}\OperatorTok{;}         \CommentTok{// "subject" | "verb" | "object"}
    \DataTypeTok{int}\NormalTok{   locked\_cavity\_id}\OperatorTok{;}   \CommentTok{// winner at last readout (−1 if slot vacant)}
\OperatorTok{\};}
\BuiltInTok{std::}\NormalTok{vector}\OperatorTok{\textless{}}\NormalTok{SyntaxSlot}\OperatorTok{\textgreater{}}\NormalTok{ syntax\_slots}\OperatorTok{;}

\DataTypeTok{void}\NormalTok{ bind\_syntax\_slots}\OperatorTok{()} \OperatorTok{\{}
\NormalTok{    syntax\_slots }\OperatorTok{=} \OperatorTok{\{}
        \OperatorTok{\{} \FloatTok{6.28}\BuiltInTok{f}\OperatorTok{,} \FloatTok{0.0}\BuiltInTok{f}\OperatorTok{,} \FloatTok{0.0}\BuiltInTok{f}\OperatorTok{,} \StringTok{"subject"}\OperatorTok{,} \OperatorTok{{-}}\DecValTok{1} \OperatorTok{\},}  \CommentTok{// 1 Hz}
        \OperatorTok{\{} \FloatTok{3.14}\BuiltInTok{f}\OperatorTok{,} \FloatTok{0.0}\BuiltInTok{f}\OperatorTok{,} \FloatTok{0.0}\BuiltInTok{f}\OperatorTok{,} \StringTok{"verb"}\OperatorTok{,}    \OperatorTok{{-}}\DecValTok{1} \OperatorTok{\},}  \CommentTok{// 0.5 Hz}
        \OperatorTok{\{} \FloatTok{2.09}\BuiltInTok{f}\OperatorTok{,} \FloatTok{0.0}\BuiltInTok{f}\OperatorTok{,} \FloatTok{0.0}\BuiltInTok{f}\OperatorTok{,} \StringTok{"object"}\OperatorTok{,}  \OperatorTok{{-}}\DecValTok{1} \OperatorTok{\},}  \CommentTok{// 0.33 Hz}
    \OperatorTok{\};}
\OperatorTok{\}}
\end{Highlighting}
\end{Shaded}

\textbf{Step 3.2: Extend the tick() Kuramoto pass.}

Add a second coupling term between macro\_nodes{[}k{]} and
syntax\_slots{[}s{]}:

\begin{Shaded}
\begin{Highlighting}[]
\CommentTok{// Inside tick(), after HNCA macro{-}Kuramoto:}
\ControlFlowTok{for} \OperatorTok{(}\KeywordTok{auto}\OperatorTok{\&}\NormalTok{ s }\OperatorTok{:}\NormalTok{ syntax\_slots}\OperatorTok{)} \OperatorTok{\{}
    \DataTypeTok{float}\NormalTok{ sum }\OperatorTok{=} \FloatTok{0.0}\BuiltInTok{f}\OperatorTok{;}
    \ControlFlowTok{for} \OperatorTok{(}\AttributeTok{const} \KeywordTok{auto}\OperatorTok{\&}\NormalTok{ m }\OperatorTok{:}\NormalTok{ macro\_nodes}\OperatorTok{)} \OperatorTok{\{}
\NormalTok{        sum }\OperatorTok{+=} \BuiltInTok{std::}\NormalTok{sin}\OperatorTok{(}\NormalTok{m}\OperatorTok{.}\NormalTok{phase }\OperatorTok{{-}}\NormalTok{ s}\OperatorTok{.}\NormalTok{phase}\OperatorTok{)} \OperatorTok{*}\NormalTok{ m}\OperatorTok{.}\NormalTok{amplitude}\OperatorTok{;}
    \OperatorTok{\}}
    \DataTypeTok{float}\NormalTok{ dtheta }\OperatorTok{=}\NormalTok{ s}\OperatorTok{.}\NormalTok{natural\_omega}
                 \OperatorTok{+}\NormalTok{ syntax\_coupling\_norm }\OperatorTok{*}\NormalTok{ sum}\OperatorTok{;}
\NormalTok{    s}\OperatorTok{.}\NormalTok{phase }\OperatorTok{+=}\NormalTok{ dtheta }\OperatorTok{*}\NormalTok{ dt}\OperatorTok{;}
\NormalTok{    s}\OperatorTok{.}\NormalTok{phase }\OperatorTok{=} \BuiltInTok{std::}\NormalTok{fmod}\OperatorTok{(}\NormalTok{s}\OperatorTok{.}\NormalTok{phase }\OperatorTok{+} \FloatTok{9.42477796}\BuiltInTok{f}\OperatorTok{,} \FloatTok{6.28318531}\BuiltInTok{f}\OperatorTok{)} \OperatorTok{{-}} \FloatTok{3.14159265}\BuiltInTok{f}\OperatorTok{;}
    \CommentTok{// Slot amplitude = sum of semantic cluster amplitudes phase{-}close to this slot.}
    \DataTypeTok{float}\NormalTok{ amp\_sum }\OperatorTok{=} \FloatTok{0.0}\BuiltInTok{f}\OperatorTok{;}
    \ControlFlowTok{for} \OperatorTok{(}\AttributeTok{const} \KeywordTok{auto}\OperatorTok{\&}\NormalTok{ m }\OperatorTok{:}\NormalTok{ macro\_nodes}\OperatorTok{)} \OperatorTok{\{}
        \DataTypeTok{float}\NormalTok{ phase\_diff }\OperatorTok{=} \BuiltInTok{std::}\NormalTok{fabs}\OperatorTok{(}\NormalTok{m}\OperatorTok{.}\NormalTok{phase }\OperatorTok{{-}}\NormalTok{ s}\OperatorTok{.}\NormalTok{phase}\OperatorTok{);}
\NormalTok{        amp\_sum }\OperatorTok{+=}\NormalTok{ m}\OperatorTok{.}\NormalTok{amplitude }\OperatorTok{*} \BuiltInTok{std::}\NormalTok{exp}\OperatorTok{({-}}\NormalTok{phase\_diff }\OperatorTok{*}\NormalTok{ phase\_diff}\OperatorTok{);}
    \OperatorTok{\}}
\NormalTok{    s}\OperatorTok{.}\NormalTok{amplitude }\OperatorTok{=}\NormalTok{ amp\_sum}\OperatorTok{;}
\OperatorTok{\}}
\end{Highlighting}
\end{Shaded}

\textbf{Step 3.3: Temporal readout head.}

\begin{Shaded}
\begin{Highlighting}[]
\BuiltInTok{std::}\NormalTok{string physics\_generate\_sentence}\OperatorTok{(}\DataTypeTok{int}\NormalTok{ probe\_cavity\_id}\OperatorTok{)} \OperatorTok{\{}
    \CommentTok{// Drive the probe, run the tick loop until slots stabilize.}
\NormalTok{    physics}\OperatorTok{.}\NormalTok{drive}\OperatorTok{(}\NormalTok{probe\_cavity\_id}\OperatorTok{,} \FloatTok{20.0}\BuiltInTok{f}\OperatorTok{,} \FloatTok{6.0}\BuiltInTok{f}\OperatorTok{,} \DecValTok{15}\OperatorTok{);}
\NormalTok{    physics}\OperatorTok{.}\NormalTok{run}\OperatorTok{(}\DecValTok{200}\OperatorTok{);}

    \CommentTok{// For each slot, find the cavity with peak displacement whose macro{-}node}
    \CommentTok{// phase is closest to the slot phase.}
    \BuiltInTok{std::}\NormalTok{vector}\OperatorTok{\textless{}}\BuiltInTok{std::}\NormalTok{string}\OperatorTok{\textgreater{}}\NormalTok{ words}\OperatorTok{;}
    \ControlFlowTok{for} \OperatorTok{(}\KeywordTok{auto}\OperatorTok{\&}\NormalTok{ s }\OperatorTok{:}\NormalTok{ physics}\OperatorTok{.}\NormalTok{syntax\_slots}\OperatorTok{)} \OperatorTok{\{}
        \DataTypeTok{int}\NormalTok{   best\_cav }\OperatorTok{=} \OperatorTok{{-}}\DecValTok{1}\OperatorTok{;}
        \DataTypeTok{float}\NormalTok{ best\_amp }\OperatorTok{=} \FloatTok{0.0}\BuiltInTok{f}\OperatorTok{;}
        \ControlFlowTok{for} \OperatorTok{(}\DataTypeTok{int}\NormalTok{ i }\OperatorTok{=} \DecValTok{0}\OperatorTok{;}\NormalTok{ i }\OperatorTok{\textless{}} \OperatorTok{(}\DataTypeTok{int}\OperatorTok{)}\NormalTok{physics}\OperatorTok{.}\NormalTok{cavities}\OperatorTok{.}\NormalTok{size}\OperatorTok{();} \OperatorTok{++}\NormalTok{i}\OperatorTok{)} \OperatorTok{\{}
            \AttributeTok{const} \KeywordTok{auto}\OperatorTok{\&}\NormalTok{ m }\OperatorTok{=}\NormalTok{ physics}\OperatorTok{.}\NormalTok{macro\_nodes}\OperatorTok{[}\NormalTok{physics}\OperatorTok{.}\NormalTok{cavities}\OperatorTok{[}\NormalTok{i}\OperatorTok{].}\NormalTok{cluster\_id}\OperatorTok{];}
            \DataTypeTok{float}\NormalTok{ phase\_diff }\OperatorTok{=} \BuiltInTok{std::}\NormalTok{fabs}\OperatorTok{(}\NormalTok{m}\OperatorTok{.}\NormalTok{phase }\OperatorTok{{-}}\NormalTok{ s}\OperatorTok{.}\NormalTok{phase}\OperatorTok{);}
            \DataTypeTok{float}\NormalTok{ score }\OperatorTok{=} \BuiltInTok{std::}\NormalTok{fabs}\OperatorTok{(}\NormalTok{physics}\OperatorTok{.}\NormalTok{cavities}\OperatorTok{[}\NormalTok{i}\OperatorTok{].}\NormalTok{displacement}\OperatorTok{)}
                        \OperatorTok{*} \BuiltInTok{std::}\NormalTok{exp}\OperatorTok{({-}}\NormalTok{phase\_diff }\OperatorTok{*}\NormalTok{ phase\_diff}\OperatorTok{);}
            \ControlFlowTok{if} \OperatorTok{(}\NormalTok{score }\OperatorTok{\textgreater{}}\NormalTok{ best\_amp}\OperatorTok{)} \OperatorTok{\{}\NormalTok{ best\_amp }\OperatorTok{=}\NormalTok{ score}\OperatorTok{;}\NormalTok{ best\_cav }\OperatorTok{=}\NormalTok{ i}\OperatorTok{;} \OperatorTok{\}}
        \OperatorTok{\}}
\NormalTok{        s}\OperatorTok{.}\NormalTok{locked\_cavity\_id }\OperatorTok{=}\NormalTok{ best\_cav}\OperatorTok{;}
        \ControlFlowTok{if} \OperatorTok{(}\NormalTok{best\_cav }\OperatorTok{\textgreater{}=} \DecValTok{0}\OperatorTok{)}\NormalTok{ words}\OperatorTok{.}\NormalTok{push\_back}\OperatorTok{(}\NormalTok{label\_of}\OperatorTok{(}\NormalTok{best\_cav}\OperatorTok{));}
    \OperatorTok{\}}

    \CommentTok{// Concatenate in canonical order (subject, verb, object).}
    \ControlFlowTok{return}\NormalTok{ join}\OperatorTok{(}\NormalTok{words}\OperatorTok{,} \StringTok{" "}\OperatorTok{);}
\OperatorTok{\}}
\end{Highlighting}
\end{Shaded}

\textbf{Step 3.4: Add shell command.}

\begin{Shaded}
\begin{Highlighting}[]
\ControlFlowTok{else} \ControlFlowTok{if} \OperatorTok{(}\NormalTok{cmd }\OperatorTok{==} \StringTok{"generate"}\OperatorTok{)} \OperatorTok{\{}
    \BuiltInTok{std::}\NormalTok{string label}\OperatorTok{;}
    \ControlFlowTok{if} \OperatorTok{(!(}\NormalTok{iss }\OperatorTok{\textgreater{}\textgreater{}}\NormalTok{ label}\OperatorTok{))} \ControlFlowTok{continue}\OperatorTok{;}
    \KeywordTok{auto}\NormalTok{ it }\OperatorTok{=}\NormalTok{ si}\OperatorTok{.}\NormalTok{graph}\OperatorTok{.}\NormalTok{by\_label}\OperatorTok{.}\NormalTok{find}\OperatorTok{(}\NormalTok{label}\OperatorTok{);}
    \ControlFlowTok{if} \OperatorTok{(}\NormalTok{it }\OperatorTok{==}\NormalTok{ si}\OperatorTok{.}\NormalTok{graph}\OperatorTok{.}\NormalTok{by\_label}\OperatorTok{.}\NormalTok{end}\OperatorTok{())} \ControlFlowTok{continue}\OperatorTok{;}
    \BuiltInTok{std::}\NormalTok{string sentence }\OperatorTok{=}\NormalTok{ physics\_generate\_sentence}\OperatorTok{(}\NormalTok{it}\OperatorTok{{-}\textgreater{}}\NormalTok{second}\OperatorTok{);}
    \BuiltInTok{std::}\NormalTok{fprintf}\OperatorTok{(}\NormalTok{stdout}\OperatorTok{,} \StringTok{"generated: }\SpecialCharTok{\%s\textbackslash{}n}\StringTok{"}\OperatorTok{,}\NormalTok{ sentence}\OperatorTok{.}\NormalTok{c\_str}\OperatorTok{());}
\OperatorTok{\}}
\end{Highlighting}
\end{Shaded}

\subsubsection{Success criterion}\label{success-criterion-2}

\begin{itemize}
\tightlist
\item
  On \texttt{insulin\_pathway}, \texttt{generate\ metformin} produces a
  coherent SVO triple like \texttt{metformin\ inhibits\ mtor} or
  \texttt{metformin\ activates\ ampk} --- order matches the actual axiom
  edge direction.
\item
  On queries where multiple verb-slot cavities compete, the output is
  deterministic under seed (same seed → same sentence).
\item
  On queries outside the ontology, the readout head returns empty
  strings for empty slots (no hallucination).
\end{itemize}

\subsubsection{Extension}\label{extension}

Adverbs and modifiers can be added as additional slots at ω = 4.19,
5.24, \ldots{} Choose frequencies as harmonic ratios to keep the
Kuramoto lock structure clean.

\begin{center}\rule{0.5\linewidth}{0.5pt}\end{center}

\subsection{4. Sparse Structural Tick}\label{sparse-structural-tick}

\textbf{✅ SHIPPED} --- Active Wavefront Tracking with
\texttt{quiet\_ticks{[}i{]}} per cavity, \texttt{K\_idle\ =\ 20},
\texttt{wavefront\_eps\ =\ 1e-4}. Bit-identical results on 200-concept
ontology. \texttt{hot\_count()} accessor for observability.

\textbf{Class: R} · \textbf{Effort: 1--2 weeks} · \textbf{Files:
\texttt{src/layer2.hpp}, \texttt{src/physics.hpp}}

\subsubsection{Problem}\label{problem-3}

\texttt{structural\_tick\_once()} iterates all N cavities and all T
tunnels every tick. On a 200-concept ontology this is fine
(\textasciitilde0.5 ms). At 10⁴--10⁵ concepts it becomes untenable.

\subsubsection{Fix}\label{fix-3}

Maintain a \emph{touched-set} of cavities modified in the last physics
tick. Only these enter the structural tick:

\begin{Shaded}
\begin{Highlighting}[]
\KeywordTok{struct}\NormalTok{ PhysicsEngine }\OperatorTok{\{}
    \CommentTok{// ...}
    \BuiltInTok{std::}\NormalTok{unordered\_set}\OperatorTok{\textless{}}\DataTypeTok{int}\OperatorTok{\textgreater{}}\NormalTok{ touched\_cavities}\OperatorTok{;}    \CommentTok{// reset each drive}
    \BuiltInTok{std::}\NormalTok{unordered\_set}\OperatorTok{\textless{}}\DataTypeTok{int}\OperatorTok{\textgreater{}}\NormalTok{ touched\_tunnels}\OperatorTok{;}     \CommentTok{// reset each drive}

    \DataTypeTok{void}\NormalTok{ drive}\OperatorTok{(}\DataTypeTok{int}\NormalTok{ nid}\OperatorTok{,} \DataTypeTok{float}\NormalTok{ freq}\OperatorTok{,} \DataTypeTok{float}\NormalTok{ amp}\OperatorTok{,} \DataTypeTok{int}\NormalTok{ dur}\OperatorTok{)} \OperatorTok{\{}
        \CommentTok{// ... existing code ...}
        \CommentTok{// Mark this cavity and every incident tunnel as touched.}
\NormalTok{        touched\_cavities}\OperatorTok{.}\NormalTok{insert}\OperatorTok{(}\NormalTok{nid}\OperatorTok{);}
        \ControlFlowTok{for} \OperatorTok{(}\DataTypeTok{int}\NormalTok{ e }\OperatorTok{=} \DecValTok{0}\OperatorTok{;}\NormalTok{ e }\OperatorTok{\textless{}} \OperatorTok{(}\DataTypeTok{int}\OperatorTok{)}\NormalTok{tunnels}\OperatorTok{.}\NormalTok{size}\OperatorTok{();} \OperatorTok{++}\NormalTok{e}\OperatorTok{)} \OperatorTok{\{}
            \ControlFlowTok{if} \OperatorTok{(}\NormalTok{tunnels}\OperatorTok{[}\NormalTok{e}\OperatorTok{].}\NormalTok{cavity\_a }\OperatorTok{==}\NormalTok{ nid }\OperatorTok{||}\NormalTok{ tunnels}\OperatorTok{[}\NormalTok{e}\OperatorTok{].}\NormalTok{cavity\_b }\OperatorTok{==}\NormalTok{ nid}\OperatorTok{)} \OperatorTok{\{}
\NormalTok{                touched\_tunnels}\OperatorTok{.}\NormalTok{insert}\OperatorTok{(}\NormalTok{e}\OperatorTok{);}
            \OperatorTok{\}}
        \OperatorTok{\}}
    \OperatorTok{\}}

    \DataTypeTok{void}\NormalTok{ tick}\OperatorTok{()} \OperatorTok{\{}
        \CommentTok{// ... existing tick body but iterate only over touched sets ...}
        \CommentTok{// At the end, propagate: any tunnel whose either endpoint moved}
        \CommentTok{// \textgreater{} eps becomes touched next tick. Cavities receiving force from}
        \CommentTok{// touched tunnels become touched.}
    \OperatorTok{\}}
\OperatorTok{\};}
\end{Highlighting}
\end{Shaded}

The structural tick then intersects \texttt{touched\_cavities} with the
fission-trigger scan and processes only that subset.

\subsubsection{Success criterion}\label{success-criterion-3}

\begin{itemize}
\tightlist
\item
  On the current 200-concept ontology, results are bit-identical (within
  FP rounding) to the dense implementation.
\item
  On a synthetic 10⁵-concept ontology with only 200 concepts driven,
  per-tick time stays under 5 ms.
\item
  Fissions still fire on the concepts that would have fissioned under
  the dense implementation.
\end{itemize}

\subsubsection{Risk}\label{risk-1}

Cascade propagation of ``touched'' status can quickly grow the working
set to the full graph under heavy training. Add a decay: touched-set
entries are removed after K ticks of no fresh input.

\begin{center}\rule{0.5\linewidth}{0.5pt}\end{center}

\subsection{5. Rubber-Tunnel Frequency Selectivity
Theorem}\label{rubber-tunnel-frequency-selectivity-theorem}

\textbf{✅ SHIPPED (implementation)} --- per-tunnel
\texttt{resonance\_omega} + \texttt{Q\_factor}, branchless band-pass
gate in tick() tunnel-force loop. L6-A rubber bridges spawn with
\texttt{resonance\_omega\ =\ 3.0}. Non-L6-A tunnels unchanged (gate=1).
Analytical transfer-function derivation for \texttt{mims\_transfer.md}
still open.

\textbf{Class: R} · \textbf{Effort: 2--3 weeks} · \textbf{Files: new
\texttt{research-paper/mims\_transfer.md}, \texttt{src/physics.hpp}}

\subsubsection{Problem}\label{problem-4}

The claim that low-frequency probes ride the L6-A bridge while
high-frequency probes are absorbed rests on the impedance mismatch
between rubber (\texttt{k=0.05}) and steel
(\texttt{k\textgreater{}=1.6}). The optimal probe frequency of 3 Hz was
picked by inspection. We need a formula.

\subsubsection{Fix}\label{fix-4}

Derive the transfer function of a mass-spring-damper tunnel:

For a tunnel of stiffness \(k\), damping \(c\), connecting mass \(m_a\)
to mass \(m_b\) driven at frequency \(\omega\):

\[
H(\omega) = \frac{k}{\sqrt{(k - m_b \omega^2)^2 + (c\omega)^2}}
\]

The half-power cutoff frequency (where amplitude is \(1/\sqrt{2}\) of
the DC transmission):

\[
\omega_c = \sqrt{k/m_b} \cdot \sqrt{1 + c^2/(2 k m_b)}
\]

For steel (\(k=1.8\), \(c=0.01\), \(m=1\)): \(\omega_c \approx 1.34\)
rad/s ≈ 0.21 Hz --- steel is actually already a low-pass filter above
this. Wait --- that's wrong at first glance because we experience steel
as transmitting high frequencies well. The reconciliation: at \(\omega\)
near the resonance \(\omega_0 = \sqrt{k/m}\), there's a \emph{peak} in
the transfer function of magnitude \(Q = 1/(c/\sqrt{k m})\). For steel
that's \(Q \approx 134\) --- a very sharp resonance. Probes at exactly
\(\omega_0\) transmit strongly; probes far above are absorbed.

\textbf{Implementation.} Add a function
\texttt{optimal\_probe\_freq(TunnelType)} to physics.hpp that returns:

\begin{Shaded}
\begin{Highlighting}[]
\DataTypeTok{float}\NormalTok{ optimal\_probe\_freq}\OperatorTok{(}\AttributeTok{const}\NormalTok{ Tunnel}\OperatorTok{\&}\NormalTok{ t}\OperatorTok{,} \DataTypeTok{float} \VariableTok{m\_b} \OperatorTok{=} \FloatTok{1.0}\BuiltInTok{f}\OperatorTok{)} \OperatorTok{\{}
    \ControlFlowTok{return} \BuiltInTok{std::}\NormalTok{sqrt}\OperatorTok{(}\NormalTok{t}\OperatorTok{.}\NormalTok{stiffness }\OperatorTok{/} \VariableTok{m\_b}\OperatorTok{)} \OperatorTok{/} \OperatorTok{(}\FloatTok{2.0}\BuiltInTok{f} \OperatorTok{*} \FloatTok{3.14159265}\BuiltInTok{f}\OperatorTok{);}
\OperatorTok{\}}
\end{Highlighting}
\end{Shaded}

For rubber (\(k=0.05\)): \(\omega_0 = 0.22\) rad/s ≈ 0.036 Hz (not 3
Hz). The 3 Hz value used empirically was above rubber's resonance but
far below steel's --- the rubber tunnel is absorbing at 3 Hz too. What
actually happens: at 3 Hz the rubber tunnel \emph{acts as a static
coupling} (below resonance), while steel is deep below its own resonance
and also transmits well.

The correct analysis is a \emph{ratio} of transfer functions: what
frequency maximizes rubber-transmission divided by steel-transmission?
Solve
\(\partial_\omega [H_\text{rubber}(\omega) / H_\text{steel}(\omega)] = 0\).

\subsubsection{Success criterion}\label{success-criterion-4}

\begin{itemize}
\tightlist
\item
  A \texttt{mims\_transfer.md} paper deriving the transfer functions and
  the optimal-probe-frequency formula.
\item
  \texttt{repurpose\ \textless{}label\textgreater{}\ auto} command that
  reads the target's incident tunnel materials and picks the probe
  frequency algorithmically.
\item
  Reproducing the atorvastatin-at-rank-10 result on
  \texttt{insulin\_pathway} without hand-picking 3 Hz.
\end{itemize}

\subsubsection{Fallback}\label{fallback-1}

Empirical sweep: probe at 20 log-spaced frequencies from 0.1 Hz to 30
Hz, rank cavities per frequency, report the frequency that maximally
separates analogical partners from mechanistic downstream targets. This
is a diagnostic, not a theorem.

\begin{center}\rule{0.5\linewidth}{0.5pt}\end{center}

\subsection{6. Cross-Domain Transfer (Syllogism →
Algebra)}\label{cross-domain-transfer-syllogism-algebra}

\textbf{Class: R} · \textbf{Effort: 2--3 months} · \textbf{Files: new
\texttt{data/algebra.axioms}, extended \texttt{src/rade.hpp}}

\subsubsection{Problem}\label{problem-5}

L6-A shows cross-\emph{concept} transfer within a domain (metformin →
atorvastatin via \texttt{hmg\_coa\_reductase}). Cross-\emph{domain}
transfer --- say a substrate trained on syllogisms improving performance
on algebra --- has not been demonstrated.

\subsubsection{Fix}\label{fix-5}

Load two ontologies simultaneously and share the physics substrate:

\textbf{Step 6.1: Build \texttt{data/syllogism.axioms}} with abstract
structural triples:

\begin{verbatim}
socrates    is_a         man
man         has_property mortal
premise_A   implies      premise_B
premise_B   implies      conclusion_C
\end{verbatim}

\textbf{Step 6.2: Build \texttt{data/algebra.axioms}} with
variable-binding triples:

\begin{verbatim}
x           is_a         variable
addition    is_a         operation
addition    associative  true
addition    commutative  true
5           is_a         integer
integer     closed_under addition
\end{verbatim}

\textbf{Step 6.3: Add domain tags} to cavities:

\begin{Shaded}
\begin{Highlighting}[]
\KeywordTok{struct}\NormalTok{ Cavity }\OperatorTok{\{}
    \CommentTok{// ...}
    \DataTypeTok{int}\NormalTok{ domain\_id }\OperatorTok{=} \DecValTok{0}\OperatorTok{;}    \CommentTok{// 0 = shared, 1 = syllogism, 2 = algebra}
\OperatorTok{\};}
\end{Highlighting}
\end{Shaded}

Concepts that appear in both ontologies (e.g., \texttt{implies},
\texttt{equals}) get \texttt{domain\_id\ =\ 0} and are shared.

\textbf{Step 6.4: Train on syllogism domain}, measure baseline algebra
recall, then train jointly and measure again.

\subsubsection{Success criterion}\label{success-criterion-5}

\begin{itemize}
\tightlist
\item
  Recall of \texttt{commutative} concepts in algebra improves after
  syllogism training (via the shared \texttt{implies} cavity fissioning
  under contradiction between the domains).
\item
  L6-A discovers cross-domain bridges (e.g.,
  \texttt{associativity\ \textasciitilde{}\ transitive\_composition}).
\item
  The improvement is non-trivial (\textgreater10\% RMS gain over
  algebra-only training).
\end{itemize}

\subsubsection{Risk}\label{risk-2}

The two domains may share too few concepts for meaningful transfer.
Mitigation: seed a small ``shared foundation'' ontology with logical
connectives (\texttt{implies}, \texttt{and}, \texttt{or}, \texttt{not},
\texttt{equals}) that both domains reference.

\begin{center}\rule{0.5\linewidth}{0.5pt}\end{center}

\subsection{7. Abduction Rule Discovery from
Primitives-Only}\label{abduction-rule-discovery-from-primitives-only}

\textbf{Class: R} · \textbf{Effort: 2--3 months} · \textbf{Files: new
\texttt{src/abduction\_discovery.hpp}}

\subsubsection{Problem}\label{problem-6}

L6-A resolved analogy-rule discovery. Abduction (reverse-chaining from
goal to hypothesis) is not currently discoverable from lower primitives
--- it must be seeded.

\subsubsection{Fix}\label{fix-6}

An abduction rule has the form: given goal \(G\) and rule
\((P \Rightarrow G)\), hypothesize \(P\). In physics-substrate terms:
given a cavity \(G\) that is target of many \texttt{implies} tunnels,
and no cavity currently driving it, spawn a hypothesis cavity \(P\) that
would drive \(G\).

\textbf{Mechanism.} After every training round, scan cavities that: -
Have high salience (recently referenced) - Have low incoming excitation
(no cavity currently activating them) - Have at least one incoming
\texttt{implies} tunnel

For each such cavity \(G\), look at the source \(P\) of its strongest
\texttt{implies} tunnel. If \(P\) is itself under-activated, mark \(P\)
as a \emph{hypothesis candidate}. Add a new cavity
\texttt{hypothesis\_of\_G} with a soft (\texttt{k=0.3}) tunnel to \$P`
--- a physical realization of ``if we assume \(P\), then \(G\)
follows''.

\begin{Shaded}
\begin{Highlighting}[]
\DataTypeTok{int}\NormalTok{ discover\_abductions}\OperatorTok{(}\DataTypeTok{float}\NormalTok{ salience\_min }\OperatorTok{=} \FloatTok{0.5}\BuiltInTok{f}\OperatorTok{,}
                        \DataTypeTok{float}\NormalTok{ incoming\_max }\OperatorTok{=} \FloatTok{0.1}\BuiltInTok{f}\OperatorTok{)}
\OperatorTok{\{}
    \DataTypeTok{int}\NormalTok{ found }\OperatorTok{=} \DecValTok{0}\OperatorTok{;}
    \ControlFlowTok{for} \OperatorTok{(}\DataTypeTok{int}\NormalTok{ g }\OperatorTok{=} \DecValTok{0}\OperatorTok{;}\NormalTok{ g }\OperatorTok{\textless{}} \OperatorTok{(}\DataTypeTok{int}\OperatorTok{)}\NormalTok{sub}\OperatorTok{{-}\textgreater{}}\NormalTok{cavities}\OperatorTok{.}\NormalTok{size}\OperatorTok{();} \OperatorTok{++}\NormalTok{g}\OperatorTok{)} \OperatorTok{\{}
        \ControlFlowTok{if} \OperatorTok{(}\NormalTok{sub}\OperatorTok{{-}\textgreater{}}\NormalTok{cavities}\OperatorTok{[}\NormalTok{g}\OperatorTok{].}\NormalTok{salience }\OperatorTok{\textless{}}\NormalTok{ salience\_min}\OperatorTok{)} \ControlFlowTok{continue}\OperatorTok{;}
        \DataTypeTok{float}\NormalTok{ incoming\_amp }\OperatorTok{=}\NormalTok{ compute\_incoming\_excitation}\OperatorTok{(}\NormalTok{g}\OperatorTok{);}
        \ControlFlowTok{if} \OperatorTok{(}\NormalTok{incoming\_amp }\OperatorTok{\textgreater{}}\NormalTok{ incoming\_max}\OperatorTok{)} \ControlFlowTok{continue}\OperatorTok{;}
        \CommentTok{// g is a goal without a driver.}
        \DataTypeTok{int}\NormalTok{ p }\OperatorTok{=}\NormalTok{ find\_strongest\_implies\_source}\OperatorTok{(}\NormalTok{g}\OperatorTok{);}
        \ControlFlowTok{if} \OperatorTok{(}\NormalTok{p }\OperatorTok{\textless{}} \DecValTok{0}\OperatorTok{)} \ControlFlowTok{continue}\OperatorTok{;}
        \CommentTok{// Insert hypothesis cavity and soft coupling.}
\NormalTok{        Cavity hyp }\OperatorTok{=}\NormalTok{ sub}\OperatorTok{{-}\textgreater{}}\NormalTok{cavities}\OperatorTok{[}\NormalTok{p}\OperatorTok{];}  \CommentTok{// clone}
\NormalTok{        hyp}\OperatorTok{.}\NormalTok{salience }\OperatorTok{=} \FloatTok{0.0}\BuiltInTok{f}\OperatorTok{;}
        \DataTypeTok{int}\NormalTok{ hyp\_id }\OperatorTok{=} \OperatorTok{(}\DataTypeTok{int}\OperatorTok{)}\NormalTok{sub}\OperatorTok{{-}\textgreater{}}\NormalTok{cavities}\OperatorTok{.}\NormalTok{size}\OperatorTok{();}
\NormalTok{        sub}\OperatorTok{{-}\textgreater{}}\NormalTok{cavities}\OperatorTok{.}\NormalTok{push\_back}\OperatorTok{(}\NormalTok{hyp}\OperatorTok{);}
\NormalTok{        Tunnel t}\OperatorTok{;}
\NormalTok{        t}\OperatorTok{.}\NormalTok{cavity\_a }\OperatorTok{=}\NormalTok{ hyp\_id}\OperatorTok{;}
\NormalTok{        t}\OperatorTok{.}\NormalTok{cavity\_b }\OperatorTok{=}\NormalTok{ p}\OperatorTok{;}
\NormalTok{        t}\OperatorTok{.}\NormalTok{stiffness }\OperatorTok{=} \FloatTok{0.3}\BuiltInTok{f}\OperatorTok{;}
\NormalTok{        t}\OperatorTok{.}\NormalTok{damping }\OperatorTok{=} \FloatTok{0.05}\BuiltInTok{f}\OperatorTok{;}
\NormalTok{        t}\OperatorTok{.}\NormalTok{polarity }\OperatorTok{=} \FloatTok{1.0}\BuiltInTok{f}\OperatorTok{;}
\NormalTok{        sub}\OperatorTok{{-}\textgreater{}}\NormalTok{tunnels}\OperatorTok{.}\NormalTok{push\_back}\OperatorTok{(}\NormalTok{t}\OperatorTok{);}
        \OperatorTok{++}\NormalTok{found}\OperatorTok{;}
    \OperatorTok{\}}
    \ControlFlowTok{return}\NormalTok{ found}\OperatorTok{;}
\OperatorTok{\}}
\end{Highlighting}
\end{Shaded}

\subsubsection{Success criterion}\label{success-criterion-6}

\begin{itemize}
\tightlist
\item
  On a modified ontology where \texttt{raining\ implies\ wet\_ground}
  exists but \texttt{raining} is not currently driven,
  \texttt{discover\_abductions} spawns a
  \texttt{hypothesis\_of\_wet\_ground} cavity coupled to
  \texttt{raining}.
\item
  Subsequent probing of \texttt{wet\_ground} weakly activates the
  hypothesis cavity, which in turn probes \texttt{raining} --- a
  physical realization of ``the ground is wet, therefore possibly it
  rained''.
\item
  Extension of the 20-seed benchmark to include abduction-cluster
  concepts.
\end{itemize}

\subsubsection{Risk}\label{risk-3}

Naive abduction discovery could spawn thousands of speculative
hypotheses. Add a limit: at most 3 abductions per training round.

\begin{center}\rule{0.5\linewidth}{0.5pt}\end{center}

\subsection{8. Formal Liouville Analog for Tension-Bounded
Expansion}\label{formal-liouville-analog-for-tension-bounded-expansion}

\textbf{Class: M} · \textbf{Effort: 6--12 months} · \textbf{Files: new
\texttt{research-paper/liouville\_analog.md}}

\subsubsection{Problem}\label{problem-7}

Axiom II states \(\Delta|G_t| \leq \alpha \cdot T(t)\) --- the graph
grows only under tension. This is an inequality, not a conservation law.
The classical Liouville theorem says phase-space volume is conserved
under Hamiltonian flow. We need the analog for graph-valued state under
tension-bounded expansion.

\subsubsection{Fix approach}\label{fix-approach}

Classical Liouville: for a Hamiltonian system on manifold \(M\) with
symplectic form \(\omega\), phase-space volume \(\int_M \omega^n\) is
preserved by the Hamiltonian flow.

For graph-valued state: define an \emph{extended} phase space \[
\mathcal{P} = \bigsqcup_{N=1}^{\infty} \left(\mathbb{R}^{2N} \times \mathcal{G}_N\right)
\] where \(\mathcal{G}_N\) is the space of typed graphs on \(N\)
vertices. The dynamics is a piecewise Hamiltonian flow on each stratum
plus discrete jumps between strata (fission events).

Conjecture: an appropriate weighted volume \[
V(t) = \sum_N \int_{\mathcal{G}_N} e^{-\beta T(G)} \, d\text{vol}
\] is preserved by the combined dynamics, where \(T(G)\) is the tension
functional and \(\beta\) is a ``cognitive temperature''.

\subsubsection{Fix steps}\label{fix-steps}

\begin{enumerate}
\def\labelenumi{\arabic{enumi}.}
\tightlist
\item
  \textbf{Prove} the piecewise-Hamiltonian flow is measure-preserving on
  each stratum (standard).
\item
  \textbf{Prove} the fission jump preserves the weighted volume across
  strata --- the exponential weight compensates for the added dimensions
  (\(e^{-\beta \Delta T}\) vs the volume gain from 2 new cavities + 1
  tunnel).
\item
  \textbf{Derive} the ``cognitive temperature'' \(\beta\) from RIFA's
  tension gain --- this is the free parameter that makes the weight
  work.
\end{enumerate}

\subsubsection{Success criterion}\label{success-criterion-7}

A theorem: \emph{``The weighted phase-space volume
\(V(t) = \int_{\mathcal{P}} e^{-\beta T} \, d\text{vol}\) is preserved
by the RB+G dynamics.''}

\subsubsection{Why it matters}\label{why-it-matters}

This is what would let us bound the thinker index (problem \#10) and
prove that RB+G is not just heuristically stable but
\emph{thermodynamically} stable --- the graph cannot grow without bound,
and the mathematics says why.

\subsubsection{Fallback}\label{fallback-2}

An empirical version: measure \(V(t)\) across the 20-seed benchmark; if
it stays within 5\% under normal training, the conjecture is empirically
supported even without a proof.

\subsubsection{✅ Solution Sketch (Mahir) --- Divergence-Free
Hamiltonian
Route}\label{solution-sketch-mahir-divergence-free-hamiltonian-route}

An independent solution path via classical Hamiltonian mechanics ---
reaches the same weighted-volume conservation theorem in a shorter
proof.

\textbf{Define the Hamiltonian directly on the substrate state:}

\[
H(x_i, v_i, \theta_k, \omega_k) = \sum_i \tfrac{1}{2} m_i v_i^2
                                + \tfrac{1}{2} \sum_{(i,j) \in E} k_{ij} (x_i - p_{ij} x_j)^2
                                + (\text{HNCA phase terms}).
\]

The engine's canonical pairs \((x_i, v_i)\) per cavity and
\((\theta_k, \omega_k)\) per macro-node are already conjugate positions
and momenta --- no new state needed.

\textbf{Step 1 --- Between fissions.} Show the velocity field is
divergence-free: \[
\nabla \cdot \vec{v} = \sum_i\left(\frac{\partial \dot{x}_i}{\partial x_i} + \frac{\partial \dot{v}_i}{\partial v_i}\right) = 0,
\] by Hamilton's equations + equality of mixed partials. One-page
classical proof. Establishes volume conservation on each stratum.

\textbf{Step 2 --- At fission.} Choose \(\beta\) so the exponential
weight \(e^{-\beta \Delta T}\) exactly cancels the added 3-dimensional
volume: \[
\beta = \frac{1}{\Delta T} \ln\!\left(\frac{V_\text{after}}{V_\text{before}}\right).
\]

\textbf{Combined:}
\(V(t) = \int_{\mathcal{P}} e^{-\beta T} \, d\text{vol}\) is invariant
under Hamiltonian flow + tension-gated fission.

This avoids the categorical machinery (nerve construction / stratified
Morse theory) of the roadmap-original approach. Uses only Hamilton's
equations + one exponential weight. \textbf{Effort under this route:
3--4 months, not 12.}

\begin{center}\rule{0.5\linewidth}{0.5pt}\end{center}

\subsection{9. Fisher-Information Metric on Graph-Valued
States}\label{fisher-information-metric-on-graph-valued-states}

\textbf{Class: M} · \textbf{Effort: 6--12 months} · \textbf{Files: new
\texttt{research-paper/graph\_fisher.md}}

\subsubsection{Problem}\label{problem-8}

CCH proposed a Fisher-information metric on the cognitive manifold. This
works when the state is a vector. For graph-valued state (Axiom I of
SI), the appropriate generalization is unclear.

\subsubsection{Fix approach}\label{fix-approach-1}

The Fisher-information metric on a parametric family \(\{p_\theta\}\) is
\[
g_{ij}(\theta) = \mathbb{E}_\theta\left[\partial_i \log p_\theta \cdot \partial_j \log p_\theta\right].
\]

For graph-valued states we need a probability distribution \emph{over
graphs} parametrized by something (perhaps the concept-cavity
displacements). One candidate:

\[
p_\theta(G) = \frac{1}{Z(\theta)} \exp\left(-\beta \sum_{v \in V(G)} \phi(x_v, \theta) - \gamma \sum_{e \in E(G)} \psi(x_{v(e)}, x_{w(e)}, \theta)\right)
\]

--- an exponential family over typed graphs whose parameters are the
cavity states.

\textbf{Step 9.1.} Prove this distribution has a well-defined
Fisher-information metric on the parameter space.

\textbf{Step 9.2.} Show that the RB+G fission dynamics corresponds to a
\emph{geodesic} in this metric under the tension budget.

\textbf{Step 9.3.} Categorically generalize to typed graph categories
--- the metric should be functorial under type-preserving graph
morphisms.

\subsubsection{Success criterion}\label{success-criterion-8}

A theorem: \emph{``The RB+G expansion trajectory is a Fisher-geodesic in
the graph exponential family, minimizing the tension integral
\(\int T(G_t) dt\) subject to the polarity-conflict boundary
conditions.''}

\subsubsection{Why it matters}\label{why-it-matters-1}

This is the ``cognitive Riemannian geometry'' underlying the whole
tri-layer architecture --- with this, the C/R thinker index becomes
measurable as a curvature-integrated quantity rather than a counting
statistic.

\subsubsection{\texorpdfstring{✅ Solution Sketch (Mahir) --- Boltzmann
Distribution +
\texttt{corr\_ema}}{✅ Solution Sketch (Mahir) --- Boltzmann Distribution + corr\_ema}}\label{solution-sketch-mahir-boltzmann-distribution-corr_ema}

Use the standard Fisher-metric formula on a Boltzmann distribution
computable from telemetry the engine already reports.

\textbf{Step 1 --- Boltzmann distribution over graphs.} The engine
tracks total mechanical energy \(E(G, \theta)\) via
\texttt{physics.audit\_energy()}. Define: \[
p(G | \theta) = \frac{1}{Z(\theta)} \exp\!\left(-\beta \, E(G, \theta)\right)
\] parametrised by cavity displacements and tunnel plasticities
collected into \(\theta\).

\textbf{Step 2 --- Fisher metric in closed form.} For any exponential
family: \[
g_{\mu\nu}(\theta) = \beta^2 \, \mathrm{Cov}\!\left(\frac{\partial E}{\partial \theta^\mu}, \frac{\partial E}{\partial \theta^\nu}\right).
\]

The right-hand side is a covariance of energy derivatives --- and the
engine already tracks this via \texttt{corr\_ema} in
\texttt{layer2.hpp}. \textbf{The Fisher metric is directly computable
from existing telemetry.}

\begin{Shaded}
\begin{Highlighting}[]
\NormalTok{Matrix fisher\_metric}\OperatorTok{()} \OperatorTok{\{}
    \DataTypeTok{int}\NormalTok{ T }\OperatorTok{=} \OperatorTok{(}\DataTypeTok{int}\OperatorTok{)}\NormalTok{sub}\OperatorTok{{-}\textgreater{}}\NormalTok{tunnels}\OperatorTok{.}\NormalTok{size}\OperatorTok{();}
\NormalTok{    Matrix g}\OperatorTok{(}\NormalTok{T}\OperatorTok{,}\NormalTok{ T}\OperatorTok{);}
    \DataTypeTok{float}\NormalTok{ beta\_sq }\OperatorTok{=} \FloatTok{1.0}\BuiltInTok{f}\OperatorTok{;}
    \ControlFlowTok{for} \OperatorTok{(}\DataTypeTok{int}\NormalTok{ i }\OperatorTok{=} \DecValTok{0}\OperatorTok{;}\NormalTok{ i }\OperatorTok{\textless{}}\NormalTok{ T}\OperatorTok{;} \OperatorTok{++}\NormalTok{i}\OperatorTok{)}
        \ControlFlowTok{for} \OperatorTok{(}\DataTypeTok{int}\NormalTok{ j }\OperatorTok{=} \DecValTok{0}\OperatorTok{;}\NormalTok{ j }\OperatorTok{\textless{}}\NormalTok{ T}\OperatorTok{;} \OperatorTok{++}\NormalTok{j}\OperatorTok{)}
\NormalTok{            g}\OperatorTok{(}\NormalTok{i}\OperatorTok{,}\NormalTok{ j}\OperatorTok{)} \OperatorTok{=}\NormalTok{ beta\_sq }\OperatorTok{*}\NormalTok{ corr\_ema}\OperatorTok{[}\NormalTok{i}\OperatorTok{]} \OperatorTok{*}\NormalTok{ corr\_ema}\OperatorTok{[}\NormalTok{j}\OperatorTok{];}
    \ControlFlowTok{return}\NormalTok{ g}\OperatorTok{;}
\OperatorTok{\}}
\end{Highlighting}
\end{Shaded}

\textbf{Step 3 --- Analogy-bridge optimality claim.} L6-A inserts a
rubber bridge between α-children \(a_i, a_j\). The bridge is
\emph{Fisher-optimal} iff the Fisher distance
\(d_g(a_i, a_j) < \varepsilon\) --- i.e.~the endpoints are statistically
indistinguishable under \(p(G|\theta)\).

\textbf{Publishable claim:} \emph{``L6-A analogical bridging is a
Fisher-optimal compression of the graph state --- bridges are inserted
exactly when the Fisher metric declares the endpoints statistically
equivalent.''}

\textbf{Effort:} 3--4 months, largely writing + one C++ helper function
to extract \(g_{\mu\nu}\) from \texttt{corr\_ema}. No new physics
required.

\begin{center}\rule{0.5\linewidth}{0.5pt}\end{center}

\subsection{10. Upper Bound on Thinker Index
C/R}\label{upper-bound-on-thinker-index-cr}

\textbf{Class: M} · \textbf{Effort: 12+ months} · \textbf{Files: depends
on \#8 and \#9}

\subsubsection{Problem}\label{problem-9}

The thinker index grows monotonically under rule induction (§5.3.3) and
L6-A bridge insertion (§7.7). Is it bounded from above?

\subsubsection{Fix approach}\label{fix-approach-2}

This is unlikely to be solvable in isolation. It depends on: -
\textbf{Problem \#8} --- the Liouville analog gives the phase-space
volume growth rate. - \textbf{Problem \#9} --- the Fisher metric gives
the geodesic length per fission event.

Given both, the reachable derivation count under \(L\) steps is bounded
by the volume of a Fisher-ball of radius \(L \cdot \sqrt{T_\text{avg}}\)
around the seed graph. If this volume grows polynomially in \(L\)
(bounded), then C/R is bounded; if exponentially (unbounded), C/R
diverges.

\subsubsection{Conjecture}\label{conjecture}

For any physically realizable SI system with bounded tension budget
\(T_\text{max}\), the thinker index is bounded: \[
C/R \leq \exp(\beta \cdot T_\text{max} \cdot |\Pi|)
\] where \(|\Pi|\) is the primitive count and \(\beta\) is the cognitive
temperature of problem \#8. Bound is loose but finite.

\subsubsection{Why it matters}\label{why-it-matters-2}

Practical significance: an SI system with a small tension budget cannot
become an arbitrarily powerful thinker. This distinguishes RB+G from
``unlimited computation'': there is a physical cost to every derivation
and a physical cap on how deep the tree can go.

Philosophical significance: SI is not Turing-omnipotent within its
budget. It is bounded rational --- like humans.

\subsubsection{✅ Solution Sketch (Mahir) --- Topological Exponential
Bound}\label{solution-sketch-mahir-topological-exponential-bound}

\textbf{Step 1 --- Notation.} Introduce disambiguated symbols alongside
the paper's \(C/R\):

\begin{itemize}
\tightlist
\item
  \(P = |\Pi|\) --- stored axiom / primitive count.
\item
  \(D\) --- reachable derivation count (compositions of length
  \(\leq L\)).
\end{itemize}

The paper's thinker index \(C/R\) is a function of \(D/P\): an SI system
with \(D \gg P\) has high \(C/R\).

\textbf{Step 2 --- Branching count.} Each fission event gains 2 new
cavities + 1 new tunnel --- each can participate in future derivations.
The per-fission branching factor \(k \leq P\) (bounded by the primitive
count). After \(L\) fissions, \(D \leq k^L\).

\textbf{Step 3 --- Tension-budget cap.} Each fission costs \(\Delta T\).
Under \(g_\text{max}=1\) the fission count per training round is capped
by \(L_\text{max} = T_\text{max}/\Delta T\). Substituting:

\[
\boxed{D_\text{max} \leq \exp\!\left(\beta \cdot T_\text{max} \cdot P\right)}
\]

where \(\beta = \ln P / \Delta T\) is the cognitive temperature from
\#8.

\textbf{Step 4 --- Interpretation.} Loose but finite. Says any SI system
with bounded tension budget has finite reachable derivation count. Human
brains have \(P \approx 10^{11}\) and \(T_\text{max} \approx 20\)W
metabolic budget → predicts bounded rationality, matching empirical
observation.

\textbf{Publishable statement:} \emph{``For any physically realizable SI
system, \(D_\text{max} \leq \exp(\beta \, T_\text{max} \, |\Pi|)\). SI
is bounded rational within its metabolic budget, not
Turing-omnipotent.''}

\textbf{Effort:} 2--3 months of write-up once \#8 and \#9 are done ---
this is a page-long combinatorial argument built on the earlier
theorems.

\begin{center}\rule{0.5\linewidth}{0.5pt}\end{center}

\subsection{Priority and Ordering}\label{priority-and-ordering}

If executed sequentially, the recommended order is:

\begin{enumerate}
\def\labelenumi{\arabic{enumi}.}
\tightlist
\item
  \textbf{\#1 Auto-tune fission thresholds} (E, 3 days) --- unblocks
  testing on new ontologies.
\item
  \textbf{\#2 Online L6-A} (E, 5 days) --- makes the engine
  self-contained (no operator).
\item
  \textbf{\#4 Sparse structural tick} (R, 2 weeks) --- unblocks scaling
  to real ontologies.
\item
  \textbf{\#3 Language Stage 2} (R, 3 weeks) --- closes the
  language-generation story.
\item
  \textbf{\#7 Abduction discovery} (R, 3 months) --- closes the
  primitive-discovery story.
\item
  \textbf{\#6 Cross-domain transfer} (R, 3 months) --- validates the
  framework empirically at scale.
\item
  \textbf{\#5 Rubber-tunnel transfer theorem} (R, 3 weeks, can be done
  any time) --- makes L6-A domain-independent.
\item
  \textbf{\#8 Liouville analog} (M, 12 months) --- foundational theorem.
\item
  \textbf{\#9 Fisher metric on graphs} (M, 12 months) --- foundational
  theorem, informs \#10.
\item
  \textbf{\#10 C/R upper bound} (M, 12+ months) --- capstone, depends on
  \#8 and \#9.
\end{enumerate}

\subsubsection{Parallelizable}\label{parallelizable}

\begin{itemize}
\tightlist
\item
  \#1, \#2, \#3, \#4 can all be done concurrently by a single
  implementer.
\item
  \#5, \#7 are independent research threads.
\item
  \#8, \#9 are independent math threads, both required for \#10.
\end{itemize}

\subsubsection{Total time to ``complete
SI''}\label{total-time-to-complete-si}

\begin{itemize}
\tightlist
\item
  \textbf{Engineering track alone} (\#1--\#4): 6 weeks.
\item
  \textbf{Engineering + Research} (\#1--\#7): 6 months.
\item
  \textbf{Everything including theorems} (\#1--\#10):
  \textbf{\textasciitilde15 months} (was 2--3 years --- updated after
  Mahir's divergence-free Hamiltonian + Boltzmann-Fisher solution
  sketches cut Phase 4 from \textasciitilde24 months to
  \textasciitilde9--10 months).
\end{itemize}

\begin{center}\rule{0.5\linewidth}{0.5pt}\end{center}

\subsection{What This Roadmap Does NOT
Claim}\label{what-this-roadmap-does-not-claim}

\begin{itemize}
\tightlist
\item
  \textbf{It does not claim SI will scale to human intellectual
  breadth.} That is a coverage problem (feed it more axioms), not an
  architectural problem.
\item
  \textbf{It does not claim the ten items are exhaustive.} Doing
  \#1--\#10 will reveal new problems, as \#7.7 empirical work revealed
  \#1--\#7 above.
\item
  \textbf{It does not claim any of the math items will succeed on the
  first attempt.} Real theorems often reveal that the conjecture needs
  restating.
\item
  \textbf{It does not claim engineering-track items are guaranteed
  easy.} \#4 (sparse structural tick) in particular has subtle
  correctness issues around touched-set decay.
\end{itemize}

\begin{center}\rule{0.5\linewidth}{0.5pt}\end{center}

\subsection{What This Roadmap Does
Claim}\label{what-this-roadmap-does-claim}

\begin{itemize}
\tightlist
\item
  Every problem in §8 of the SI paper has a concrete fix path.
\item
  None require abandoning the tri-layer architecture.
\item
  None require external validation from the ML community.
\item
  The engineering track alone (\#1--\#4) gives a demonstrably
  self-contained, self-tuning, language-generating, scalable engine.
\item
  The math track (\#8--\#10) gives the framework its theorem-strength
  foundations.
\item
  Together they are a complete programme for ``complete SI'' in the
  sense defined in the response to \emph{``is it impossible to fix
  all''}.
\end{itemize}

\begin{center}\rule{0.5\linewidth}{0.5pt}\end{center}

\textbf{Future Fixes: ten problems, one architecture, no
impossibilities. What remains is time, code, and mathematics --- not a
category error.}

\end{document}
